\section{Design Principles}\label{sec:principles}

The work was conducted under a single guiding constraint: \emph{do
not seek the simplest solution; seek the solution that is most
coherent with the project, conforms to MIL-STD-882E hazard
tracking practice, maximizes resilience, and preserves performance
under deterministic bounds}. This constraint translates into four
operational axes (\cref{tab:axes}), invoked explicitly in
subsequent sections when justifying a design decision.

\begin{table*}[t]
\centering
\small
\caption{Design principles. Each axis is invoked at most once per
technical sub-section; references to the axes appear as
``per axis~3.X''.}
\label{tab:axes}
\begin{tabularx}{\textwidth}{p{2.6cm}X}
\toprule
\textbf{Axis} & \textbf{Operational meaning} \\
\midrule
3.1 Coherence
  & Single source of truth for every on-disk invariant in C, via
    \texttt{offsetof()} macros guarded by \texttt{BUILD\_BUG\_ON}
    in the kernel and \texttt{static\_assert} in user-space.
    Constants are mirrored byte-for-byte; drift is caught at
    build time. \\
\addlinespace[0.3em]
3.2 MIL-STD-882E conformance
  & Each correctable hazard is a distinct journal entry, captured
    with block number, timestamp, symbol count and CRC.
    Aggregate counters are rejected as
    forensically insufficient~\cite{milstd882e}. \\
\addlinespace[0.3em]
3.3 Resilience over throughput
  & Where tradeoffs arise, FTRFS chooses correction over
    detection-only, fail-closed over best-effort, layered defense
    over single checksum, universal coverage over opt-in. The
    deprecated \texttt{INODE\_OPT\_IN} scheme is replaced by
    unconditional \texttt{INODE\_UNIVERSAL}. \\
\addlinespace[0.3em]
3.4 Bounded performance
  & Hot paths use kernel-accelerated CRC32 and shortened RS codes
    from \texttt{lib/reed\_solomon}. The decoder runs only on
    CRC failure ($>$99\% of reads bypass it). O(1) per block, no
    FPU, no runtime division on metadata paths. The future
    Q16.16 entropy LUT (format v4) is bounded analytically at
    $\lceil N/2 \rceil$ Q-units, five orders of magnitude below
    the Family A/B classification threshold. \\
\bottomrule
\end{tabularx}
\end{table*}
